\section{Numerical illustrations}
\label{sec:numerics}

We illustrate the mechanisms established above on a small number of diagnostic
examples, each tied to a specific statement of the preceding sections. The
deterministic Gaussian and separable log-concave expectations are evaluated
analytically or by fixed Gauss--Hermite quadrature, and the nonseparable
affine-mode target uses a fixed scrambled Sobol design. The stochastic runs use
float64 arithmetic with explicit seeds, and each baseline and sticking-the-landing
pair is driven by the same random-number stream, so that the comparison is
paired. The target formulas, grids, seeds, quadrature rules, and stopping
conditions are recorded in \Cref{app:experiments}.

\subsection{The three stages}
\label{sec:numerics-three-stage}

The covariance-scaling sweep uses diagonal Gaussian targets with initial
covariances as small as $\lambda_{0,\min} = 10^{-10}$, and we record as the
burn-in endpoint the first iterate with $\lambda_{\min}(C_n) \geq 1/(2\beta)$.
Across the sweep this time is linear in $\logp (1/(\beta \lambda_{0,\min}))$,
which is the burn-in count of the two-sided covariance recursions in
\Cref{lem:cov-Riem,lem:cov-KL}. \Cref{fig:num-three-stage} shows one hard KL run
in dimension $N_\theta = 20$ with $\kappa = 1000$, where the three stages of
\Cref{thm:three-stage} are visible in sequence: the covariance calibrates, the
energy then decays rapidly, and the Fisher--Rao mean-block drift
$\| C_n^{1/2} g(a_n) \|$ decays last.

\begin{figure}[t]
\centering
\begin{subfigure}{0.47\linewidth}
  \centering
  \includegraphics[width=\linewidth]{fig_covboot_three_stage}
  \caption{A hard KL trajectory.}
\end{subfigure}\hfill
\begin{subfigure}{0.51\linewidth}
  \centering
  \includegraphics[width=\linewidth]{fig_covboot_wasserstein_bootstrap}
  \caption{One transport step, then a Fisher--Rao tail.}
\end{subfigure}
\caption{Covariance initialization and the three stages. On the left, a Gaussian
run with $N_\theta = 20$, $\kappa = 1000$, and $\lambda_{0,\min} = 10^{-10}$
separates covariance burn-in, global energy localization, and the local regime.
On the right, one Bures--Wasserstein step with $c = 1/2$, hence transport
stepsize $\eta = c/\beta$, lifts the smallest covariance eigenvalue to the
curvature scale irrespective of the initial underdispersion, and the subsequent
Fisher--Rao warm-up is essentially flat in $\lambda_{0,\min}$, as in
\Cref{thm:burnin-BW}.}
\label{fig:num-three-stage}
\end{figure}

The transport panel supports one claim only. One Bures--Wasserstein
step removes the small-covariance cost and hands the trajectory back to a
Fisher--Rao scheme; it does not change the $\kappastar$-scale obstruction of
\Cref{thm:sharp-cont}, and it is not a blended Wasserstein--Fisher--Rao flow.

\subsection{The two discretizations}
\label{sec:numerics-discretization}

The deterministic comparison uses two-dimensional Gaussian and smooth
log-concave targets together with a convex quartic target that satisfies neither half of
\Cref{assump:logconcave-smooth}, its Hessian being unbounded above and degenerate along one
direction. Both schemes use the same explicit mean
update and evaluate the same population score and Hessian. We report the
iteration count and the elapsed flow time $n \Delta t$ separately, the latter
being a discretization diagnostic rather than a timing measurement.

\begin{figure}[t]
\centering
\begin{subfigure}{0.58\linewidth}
  \centering
  \includegraphics[width=\linewidth]{fig_convergence_speed_smooth_logconcave}
  \caption{Matched-step trajectories on a smooth target.}
\end{subfigure}\hfill
\begin{subfigure}{0.39\linewidth}
  \centering
  \includegraphics[width=\linewidth]{fig_iterations_vs_stepsize}
  \caption{Iterations to tolerance.}
\end{subfigure}
\caption{The Riemannian and KL discretizations. At matched stable stepsizes the
two trajectories have comparable first-order behavior, and the Riemannian map
often contracts faster per step in the tested range, while the KL covariance
resolvent remains well defined at every stepsize and avoids a matrix
exponential. The stepsize windows of \Cref{thm:glob-Riem,thm:glob-KL} are
sufficient and conservative in these low-dimensional examples.}
\label{fig:num-discretization}
\end{figure}

The robustness recorded here is a property of the covariance map alone. On a
log-concave target one has $C_n^{-1} + \Delta t A_n \succ 0$ for every
$\Delta t > 0$, so the KL covariance update cannot leave the positive-definite
cone, which is the content of \Cref{lem:cov-KL}. The full algorithm can still
fail at a large stepsize through the explicit mean update, which carries the
factor $\Delta t\, C_n$.

\subsection{Sticking-the-landing variance and the constant-step floors}
\label{sec:numerics-stl}

The stochastic grid uses Gaussian targets with
$N_\theta \in \{2, 8, 32\}$ and $\kappa \in \{10, 100, 10^4\}$, together with
smooth $\log \cosh$ targets with $N_\theta \in \{2, 8, 32\}$ and
$\kappa \in \{10, 100\}$, each cell using twenty seeds. The energy plotted is
evaluated independently of the update sample.

\begin{figure}[t]
\centering
\begin{subfigure}{0.48\linewidth}
  \centering
  \includegraphics[width=\linewidth]{fig_stl_variance_by_distance_to_optimum}
  \caption{Variance ratio versus mean displacement.}
\end{subfigure}\hfill
\begin{subfigure}{0.49\linewidth}
  \centering
  \includegraphics[width=\linewidth]{fig_stl_batchsize_effect}
  \caption{Tail energy versus batch size.}
\end{subfigure}
\caption{Price/Hessian sticking-the-landing diagnostics. On a non-Gaussian
$\log\cosh$ target, the mean-block Fisher--Rao variance ratio is positive at the
optimizer and increases with the mean displacement in every plotted
configuration, which is the displacement-dependent part $V_1 \| a - a_\star \|^2$
of the local variance decomposition in \Cref{lem:loc-variance}. Mini-batching
reduces the median tail energy gap approximately as $1/B$, which is the batch
scaling of the Hessian-Lipschitz residual floors in
\Cref{thm:stoch-glob-Riem,thm:stoch-glob-KL}; the assumption-minimal floors of the
same theorems do not scale in $B$ at their prescribed stepsizes. The runs start from a
generic initialization at a fixed stepsize rather than from the quadratic rescue those
theorems assume, so the panels illustrate the scaling rather than instantiate the
bounds.}
\label{fig:num-stl}
\end{figure}

These panels support the variance decomposition of \Cref{sec:stochastic}. They
do not test the finite-horizon containment bound of \Cref{prop:exit-prob}, whose
horizon-dependent batch requirement is a worst-case statement and not a
prediction for the median trajectory.

\subsection{Affine-invariant modes}
\label{sec:numerics-affine}

The affine-metric experiment varies $\tau \in \{-\omega/(2N_\theta), 0,
\omega/(2N_\theta)\}$ for several values of $\omega$, with five initializations
isolating mean, volume-high, volume-low, traceless shape, and mixed
perturbations. \Cref{fig:num-affine} reports the time-to-tolerance ratio
$T(\tau)/T(0)$. At $\tau = -\omega/(2N_\theta)$ the trace rate
$1/(2(\omega + N_\theta \tau))$ of \Cref{thm:aff-modes} doubles, and the
volume-dominated runs take about half as long; at $\tau = \omega/(2N_\theta)$
their time scale grows by about $3/2$. The mean- and traceless-dominated runs
remain near one, as the corresponding modal rates $1/\eta$ and $1/(2\omega)$ do
not involve $\tau$.

\begin{figure}[t]
\centering
\begin{subfigure}{0.49\linewidth}
  \centering
  \includegraphics[width=\linewidth]{fig_ot_gaussian_tau_speedup}
  \caption{Gaussian target.}
\end{subfigure}\hfill
\begin{subfigure}{0.49\linewidth}
  \centering
  \includegraphics[width=\linewidth]{fig_ot_logconcave_tau_speedup}
  \caption{Strongly log-concave target.}
\end{subfigure}
\caption{Mode preconditioning in the affine-invariant family. The Gaussian panel
calibrates the exact modal computation of \Cref{thm:aff-modes}; the
non-Gaussian panel is an illustration of the same preconditioning effect and
not a convergence theorem for the full metric family, in the sense of
\Cref{rem:aff-scope}.}
\label{fig:num-affine}
\end{figure}

\subsection{Beyond log-concavity}
\label{sec:numerics-nonconvex}

For the double well of \Cref{sec:global}, all expectations are deterministic
Gauss--Hermite quadratures and all trajectories are centered, so that the mean
term of the split stationarity measure vanishes by symmetry and the plotted
quantity is its covariance term.

\begin{figure}[t]
\centering
\begin{subfigure}{0.32\linewidth}
  \centering
  \includegraphics[width=\linewidth]{fig_nonconvex_riemannian_cascade}
  \caption{Riemannian cascade.}
\end{subfigure}\hfill
\begin{subfigure}{0.32\linewidth}
  \centering
  \includegraphics[width=\linewidth]{fig_nonconvex_kl_pole}
  \caption{KL pole.}
\end{subfigure}\hfill
\begin{subfigure}{0.32\linewidth}
  \centering
  \includegraphics[width=\linewidth]{fig_nonconvex_clipped_kl_stationarity}
  \caption{Clipped KL stationarity.}
\end{subfigure}
\caption{The non-log-concave boundary. Although the Hessian satisfies
$\abs{V_R''} \leq 1$, negative curvature drives the covariance cascade of
\Cref{prop:cascade-Riem} and the rational pole of \Cref{prop:pole-KL}. With the
covariance constrained to $[\lambda_-, \lambda_+]$ and a stepsize in the window
of \Cref{thm:clip-KL}, the running minimum of the centered covariance KL stays
below the energy-drop envelope. Once the upper constraint binds, the iterate
remains pinned and the next displacement vanishes, which is constrained
stationarity and not a small unconstrained Fisher--Rao gradient.}
\label{fig:num-nonconvex}
\end{figure}

% SOURCES:
% codex-draft/sections/07-numerics.tex (all subsections) -> re-voiced, compressed
%   from five to five figures with the affine-mode panel retained as support for
%   thm:aff-modes; all cross-references retargeted to the new label namespace.
